% =====================================================================
% Rewrite (2026-07-09) of the learning-dynamics passage, incorporating
% the best-vs-last TRAIN-split diagnostic:
%   scripts/eval_bestcombo_best_vs_last_trainsplit.py
%   -> exp/combined_resolution/trainsplit_best_vs_last.json
% (best-val-MAE checkpoint and the epoch-1500 latest_state.pth, both
% scored on the 840 training images; val values from the logged curves).
% Corrections vs the old text: (1) best epochs are NOT all within 250;
% (2) CSRNet does not plateau — it mildly but classically overfits;
% (3) MobileCount's 3.7-3.8x rise is NOT memorisation but optimisation
% instability; (4) the capacity conclusion is hedged accordingly.
% =====================================================================

The learning curves provide insight into this behaviour. Both models
attain their best validation checkpoint within the first few hundred
epochs (between epochs $105$ and $445$ across networks and
resolutions), but their subsequent training dynamics differ sharply —
and, as scoring the end-of-training state on the \emph{training} split
reveals, in the opposite way to what the validation curves alone
suggest.
% SOURCE: best epochs — CSRNet 33/237/265/370/424/216, MobileCount
% 176/309/327/445/105/224 (49x33 -> 1544x1038); verified 2026-07-09.

CSRNet's validation MAE is not a plateau. It continues to improve
genuinely until epochs ${\sim}220$--$420$, depending on resolution,
and drifts upward by $12$--$34\%$ over the remainder of the budget.
The training-split diagnostic identifies this late drift as mild but
textbook overfitting: while validation error rises, CSRNet's training
error keeps falling to the very end of training (from $0.16$ at its
best checkpoint to $0.07$ at epoch $1500$ at $386\times260$; from
$0.09$ to $0.04$ at $1544\times1038$), finishing $5$--$8\times$ below
the validation error. The final ${\sim}1{,}100$ epochs therefore trade
generalisation for memorisation of the $840$ training images — slowly
enough that best-checkpoint selection absorbs the damage.
% SOURCE: trainsplit_best_vs_last.json — CSRNet best->last
% (train/val): 193x130 0.29/0.53 -> 0.21/0.60; 386x260 0.16/0.29 ->
% 0.07/0.32; 772x519 0.08/0.19 -> 0.04/0.23; 1544x1038 0.09/0.18 ->
% 0.04/0.24. Last-200-epoch val median / best val = 1.11-1.34.

MobileCount's far larger degradation — it finishes at
\textbf{$3.7$--$3.8\times$} its best validation error at the two
largest resolutions — is, counter-intuitively, \emph{not}
memorisation. The same diagnostic shows its training error rising in
lockstep with the validation error (train/val $0.44/0.43 \rightarrow
1.46/1.59$ at $772\times519$ and $0.33/0.43 \rightarrow 1.49/1.58$ at
$1544\times1038$): the final model is equally worse on images it saw
at every epoch, which memorisation cannot produce. Moreover, the
damage occurs early: the median validation MAE already sits near
$1.5$ by epochs $250$--$500$, when the exponential schedule still
retains ${\sim}30\%$ of the base learning rate. The optimum found at
epochs $105$--$224$ is subsequently \emph{lost} — a signature of
optimisation instability rather than of data fitting. The mechanism
cannot be isolated from these runs, but the pattern implicates the
transported training configuration rather than the data: the
optimisation hyperparameters were validated only at $386\times260$,
and stability degrades with distance from that operating point in
\emph{both} directions. Within $97\times65$--$386\times260$ the recipe
is well behaved (end-of-training degradation of $1.2$--$1.7\times$, of
the same order as CSRNet's drift); at $4$--$16\times$ the tuned pixel
count MobileCount's degradation grows to $3.7\times$; and at the
opposite extreme of $49\times33$ — where, among other regime changes,
each gradient step is estimated from ${\sim}60\times$ fewer pixels —
\emph{both} networks show the same both-splits collapse. Whether the
high-resolution failure traces to the fixed learning rate of $10^{-3}$
acting on changed gradient statistics, or to another element of the
recipe, cannot be separated without per-resolution retuning.
% SOURCE: trainsplit_best_vs_last.json — MobileCount best->last
% (train/val): 386x260 0.40/0.39 -> 0.46/0.49; 772x519 0.44/0.43 ->
% 1.46/1.59; 1544x1038 0.33/0.43 -> 1.49/1.58; mid-range degradation
% 97x65 x1.2-1.3, 193x130 x1.3-1.6. 49x33: CSRNet 5.94/5.96 ->
% 12.63/12.52 (x2.1), MobileCount 2.61/2.68 -> 6.29/6.12 (x2.3). Val
% median over ep 251-500: 1.58 (772), 1.45 (1544). LR at ep 250:
% 0.995^250 = 0.29. Pixels/step: 49x33 = 1,617 px vs 386x260 =
% 100,360 px per image (~62x).

Best-checkpoint selection shields the reported test results from both
failure modes, since the memorisation is slow and the instability
strikes only after the optimum has been recorded. It does, however,
weaken a purely capacity-based reading of MobileCount's saturation:
because MobileCount never enters a memorisation regime at any
resolution — its capacity is evidently not exhausted by the $840$
training images — and its high-resolution runs are terminated by
instability rather than by saturated learning, the present data cannot
distinguish whether the network \emph{cannot represent} the additional
spatial information or was simply \emph{not trained stably enough} to
exploit it under hyperparameters transported from $386\times260$.
CSRNet, whose training remains stable at every resolution above
$49\times33$, converts additional resolution into accuracy either way.
